% ===================================================================
%  పట్టిక సంఖ్య 4 — పరిణతి మరియు వృద్ధాప్యం.
%
%  Traduction, faite en 2026, en telougou d'aujourd'hui, sur l'ido de
%  texto/io. Regles de la colonne : voir 10-pattika-01.tex. Deux
%  scenes, deux numerotations : les cles portent c1 et c2.
%
%  LA PARENTE TELOUGOUE EST LA PLUS PRECISE DU RELEVE, ET CE TABLEAU
%  L'OBLIGE A TOUT DIRE. Le vietnamien et le marathi avaient deja
%  force la main au meme endroit ; le telougou va plus loin, parce
%  qu'il n'a PAS de mot vague a offrir :
%
%    అల్లుడు      le gendre        కోడలు      la bru
%    మామగారు      le beau-pere     అత్తగారు   la belle-mere
%    బావ          le beau-frere    ఆడబిడ్డ    la belle-sœur (sœur du mari)
%    అక్క         la sœur ainee    చెల్లెలు   la sœur cadette
%    మనవడు        le petit-fils    మనవరాలు    la petite-fille
%    ముత్తాత      l'arriere-grand-pere
%    తాతయ్య       le grand-pere    నాయనమ్మ    la grand-mere paternelle
%
%  ET LA GRAND-MERE DU MARIAGE A DU ETRE TRANCHEE. Le telougou
%  distingue నాయనమ్మ, la mere du pere, d'అమ్మమ్మ, la mere de la mere ;
%  le livret ne dit pas de quel cote vient celle du renvoi (29). Au
%  deuxieme tableau le texte tranche lui-meme — ce sont les parents
%  de Ioannes qui vieillissent —, donc నాయనమ్మ et తాతయ్య. Au
%  premier, rien ne tranche : on ecrit అవ్వ, le mot qui ne prend pas
%  parti. Choisir au hasard aurait ajoute au livret une information
%  qu'il n'a pas.
%
%  LE RESTE EST TELOUGOU SANS EFFORT : చామంతి le chrysantheme —
%  la fleur des fetes d'ici —, కౌజుపిట్ట la perdrix, పొదలు les
%  buissons, ద్రాక్ష గెలలు les grappes, ముత్యాల హారం le collier de
%  perles, కడియం le bracelet, చెవిపోగులు les boucles d'oreilles,
%  నగలు les bijoux, మంచు la neige, చంద్రుడు la lune, నక్షత్రాలు les
%  etoiles, చీకటి l'obscurite, కషాయం la tisane, మందుల చీటీ
%  l'ordonnance, పంటినొప్పి le mal de dents, కీళ్ళనొప్పులు les
%  rhumatismes, దిండు le coussin, పీట le tabouret, దినదర్శిని le
%  calendrier, చిత్రపటం le portrait.
%
%  LA ROBE BLANCHE RESTE UNE ROBE BLANCHE. La mariee de ce tableau
%  porte une robe europeenne a fleurs d'oranger ; ecrire చీర
%  donnerait une phrase plus naturelle et une noce qui n'est pas
%  celle du livret. On modernise la langue, jamais les choses — la
%  regle sert ici pour la seizieme fois.
%
%  SEPT PIEGES PRIS D'AVANCE, ET SEPT AUTRES QUI ONT PASSE QUAND
%  MEME. Ce tableau donne enfin la regle exacte, que les tableaux 1 a
%  3 n'avaient qu'approchee.
%
%  CE N'EST PAS L'AGGLUTINATION, C'EST L'ORDRE MODIFICATEUR-TETE. En
%  telougou TOUT ce qui qualifie precede ce qui est qualifie, sans
%  exception : l'adjectif, le complement de nom, la relative, le
%  participe, le groupe instrumental. Donc DES QU'UN RENVOI TOMBE SUR
%  UN MODIFICATEUR ET UN AUTRE SUR SA TETE, les deux sortent a
%  l'envers — mecaniquement, quelle que soit la phrase.
%
%  D'ou la liste des sept qui ont passe, toutes de cette forme :
%  « le manteau A COL DE FOURRURE » (8 devant 7), « l'aveugle DEVANT
%  LA PHARMACIE » (16 devant 20), « l'ordonnance QUE LE MEDECIN A
%  ECRITE » (25 devant 27), « la robe de chambre A BOUTONS DE NACRE »
%  (36 devant 35), « le grand-pere QUE LES RHUMATISMES CLOUENT AU
%  FAUTEUIL » (34 devant 33). Aucune n'est un hasard de redaction :
%  ce sont toutes des tetes precedees de leur modificateur.
%
%  LE REMEDE EST DONC MECANIQUE LUI AUSSI : quand deux renvois sont
%  en rapport modificateur-tete, poser la tete d'abord et rejeter le
%  modificateur derriere, en apposition ou en relative detachee. Ce
%  qui se verifie maintenant AVANT d'ecrire, en cherchant les paires
%  de renvois liees par « de », « avec », « qui ».
% ===================================================================

%%K t04-apar-1 apar
\VUsaut{4.0mm}
\VUtitre{15.0pt}{60}{పట్టిక సంఖ్య 4}
\VUsaut{1.6mm}
\VUtitre{11.4pt}{0}{పరిణతి మరియు వృద్ధాప్యం.}
\VUsaut{3.2mm}

%%K t04-tit-1 sub
\VUcentre{11.4pt}{0}{\textit{మొదటి రంగం.}}
\VUsaut{1.6mm}
\VUtitre{11.4pt}{0}{పెళ్ళి విందు.}
\VUsaut{2.6mm}

%%K t04-tit-2 sub
\VUpk{10.2pt}{40}{శరదృతువు.}
\VUsaut{3.0mm}

%%K t04-c1-01-1 p
1. --- యువకులు పెద్దవాళ్ళయ్యారు. వాళ్ళలో ఒకడు, జాన్, పెళ్ళి చేసుకుంటున్నాడు. కొత్త దంపతుల
కుటుంబాలను ఒకే బల్ల చుట్టూ చేర్చే \VUgras{పెళ్ళి విందుకు} మనం హాజరవుతున్నాం.

%%K t04-c1-02-1 p
2. --- మనం \VUgras{శరదృతువులో} ఉన్నాం, \VUgras{సెప్టెంబరు నెలలో}. \VUgras{అక్టోబరు},
\VUgras{నవంబరు} నెలల చలిగాలి ఇంకా పల్లె పచ్చని ఆకులను దోచుకోలేదు. సాయంత్రపు గాలి వెచ్చగా,
సువాసనతో ఉంది; అందుకే భోజనం తోట పక్కనున్న \VUgras{వరండా} \textsuperscript{(8)} కింద
జరుగుతోంది.

%%K t04-c1-03-1 p
3. --- దూరంగా \VUgras{కొండపై} \textsuperscript{(3)} జాన్ తల్లిదండ్రుల ఇల్లు ఉంది — పచ్చదనంలో
దాగిన హుందా అయిన \VUgras{భవనం} \textsuperscript{(1)}; ఉత్తమమైన ద్రాక్షసారా ఇచ్చే అందమైన
ద్రాక్షతోటలు ఆ కొండ వాలును కప్పుతున్నాయి. ద్రాక్షతోట రైతు వాటిని సాగుచేసి కాపాడుతాడు.

%%K t04-c1-04-1 p
4. --- ద్రాక్షతోటల మీదుగా \VUgras{కౌజుపిట్టల} \textsuperscript{(2)} గుంపు ఎగురుతోంది,
\VUgras{వేటకాపరి} \textsuperscript{(5)} దగ్గరికి రావడంతో అవి భయపడ్డాయి. అతను చంకలో
\VUgras{తుపాకితో}, భుజాన \VUgras{వేట సంచితో}, \VUgras{పొదలను} \textsuperscript{(4)}
శోధిస్తున్నాడు, తన \VUgras{కుక్కతో} \textsuperscript{(6)}. ఆ ఆస్తిని అతను కనిపెడతాడు,
చోరవేటగాళ్ళు వేటజంతువులను నాశనం చేయకుండా అడ్డుకుంటాడు.

%%K t04-c1-05-1 p
5. --- \VUgras{వరండాకు} బయట ఒక \VUgras{ద్రాక్ష పందిరి} ఎక్కుతోంది, దాని నుంచి దట్టమైన
\VUgras{ద్రాక్ష గెలలు} \textsuperscript{(7)} వేలాడుతున్నాయి.

%%K t04-c1-06-1 p
6. --- \VUgras{బల్లను} \textsuperscript{(16)} అలంకరించే అందమైన శరదృతువు పూలు \VUgras{చామంతులు}
\textsuperscript{(9)}; పెళ్ళికూతురి \VUgras{తండ్రి} \textsuperscript{(10)} వాటిలో ఒకదాన్ని తన
ఫ్రాక్ కోటు గుండీ రంధ్రంలో పెట్టుకున్నాడు.

%%K t04-c1-07-1 p
7. --- భోజనం ముగుస్తోంది. \VUgras{నౌకరు} \textsuperscript{(11)} తీపి వంటకాలు తేవడం
మొదలుపెట్టాడు, ఇక అవసరం లేని \VUgras{చెంచాలు} \textsuperscript{(14)}, \VUgras{ఫోర్కులు}
\textsuperscript{(12)}, \VUgras{కత్తులు} \textsuperscript{(13)}, \VUgras{పెద్ద పళ్ళేలు}
\textsuperscript{(15)} — భోజన సామగ్రిలోని ఈ రకరకాల వస్తువులను త్వరలో తీసేస్తాడు.

%%K t04-tit-3 sub
\VUpk{10.2pt}{40}{బంధువులు.}
\VUsaut{3.0mm}

%%K t04-c1-08-1 p
8. --- అందరూ ఆనందంగా కనిపిస్తున్నారు, పెళ్ళికొడుకు \VUgras{తల్లి} \textsuperscript{(17)} తన
పిల్లలను చూసే చిరునవ్వు ఆమె సంతోషాన్ని చాటుతోంది.

%%K t04-c1-09-1 p
9. --- ఈ పెళ్ళి ఆ ప్రాంతంలోని రెండు ధనిక కుటుంబాలను కలుపుతోంది. \VUgras{పెళ్ళికొడుకు}
\textsuperscript{(18)} జాన్ ఒక పెద్ద భూస్వామి \VUgras{కొడుకు}, ఇప్పుడు అతను నేర్పరి అయిన ఒక
పరిశ్రమాధిపతికి \VUgras{అల్లుడు} అవుతున్నాడు.

%%K t04-c1-10-1 p
10. --- \VUgras{దంపతులు} చిన్నవాళ్ళే, అయినా చాలాకాలం నుంచి వాళ్ళ \VUgras{నిశ్చితార్థం}
జరిగింది. \VUgras{యువతి} \textsuperscript{(19)} మార్తా, నారింజ పూలతో అలంకరించిన తన
\VUgras{తెల్లని పెళ్ళి దుస్తుల్లో} అందంగా ఉంది.

%%K t04-c1-11-1 p
11. --- ఆమె \VUgras{మామగారు} \textsuperscript{(20)} అంత మృదువైన \VUgras{కోడలు} దొరికినందుకు
చాలా సంతోషంగా ఉన్నారు, జాన్ \VUgras{అత్తగారు} \textsuperscript{(21)} తన \VUgras{కూతురు} అంత
అందంగా ఉండటం చూసి చిరునవ్వు నవ్వుతున్నారు.

%%K t04-c1-12-1 p
12. --- బల్ల చివర మిగతా \VUgras{అతిథులు} \textsuperscript{(22)} కూర్చున్నారు —
\VUgras{బంధువులు, స్నేహితులు, చిన్నాన్న కొడుకులు, చిన్నాన్న కూతుళ్ళు}. చంకలో సర్వ్ చేసే
తువ్వాలుతో \VUgras{వడ్డన పర్యవేక్షకుడు} \textsuperscript{(23)} శ్రద్ధగా వాళ్ళకు
వడ్డిస్తున్నాడు.

%%K t04-tit-4 sub
\VUpk{10.2pt}{40}{స్నేహితులు.}
\VUsaut{3.0mm}

%%K t04-c1-13-1 p
13. --- బల్లకు కుడివైపు ఇదిగో ఒక \VUgras{కుమారి} \textsuperscript{(24)}, పెళ్ళికూతురి
\VUgras{స్నేహితురాలు}; ఆ తర్వాత ఆమె \VUgras{అన్న}, అతనే ఆమె \VUgras{గౌరవ సహచరుడు}
\textsuperscript{(25)}. అతను తన \VUgras{గౌరవ కుమారితో} \textsuperscript{(26)} కబుర్లు
చెబుతున్నాడు — ఆమె పెళ్ళికొడుకు చెల్లెలు, ఈ పెళ్ళితో మార్తాకు \VUgras{ఆడబిడ్డ} అవుతుంది. ఆమె
అందమైన యువతి. ఆమెకు మంచి \VUgras{నగలు}, \VUgras{ముత్యాల హారం}, బంగారు \VUgras{కడియం} ఉన్నాయి.

%%K t04-c1-14-1 p
14. --- వాళ్ళ దగ్గర నిలబడి తన గ్లాసు ఎత్తుతున్న పెద్దమనిషి కుటుంబ \VUgras{మిత్రుడు}
\textsuperscript{(27)}. అతను పెళ్ళికొడుకు \VUgras{సాక్షుల్లో} ఒకడు. అతను
\VUgras{శుభాకాంక్షలు తెలుపుతూ} కొత్త దంపతుల \VUgras{ఆరోగ్యానికి} తాగుతున్నాడు, వాళ్ళకు
దీర్ఘకాల సౌభాగ్యం కోరుతున్నాడు. అతను వేడుక దుస్తుల్లో ఉన్నాడు — ఫ్రాక్ కోటు,
\VUgras{పాలిష్ బూట్లు} \textsuperscript{(28)}. అతను \VUgras{అవ్వకు} \textsuperscript{(29)}
తోడుగా వచ్చినవాడు, ఆమె \VUgras{వితంతువు}.

%%K t04-c1-15-1 p
15. --- \VUgras{ఫ్రాక్ కోటులో} \textsuperscript{(31)}, సన్నని నల్ల మీసంతో ఉన్న యువకుడు జాన్
\VUgras{బావ} \textsuperscript{(30)}. అతను ప్రసిద్ధ ఇంజనీరు. చాలా హుందా అయిన ఒక \VUgras{యువతి}
\textsuperscript{(33)} పక్కన అతను ఉన్నాడు — ఆమె మార్తా \VUgras{అక్క}, తన చిన్నాన్న కొడుకు చేయి
పట్టుకుని పువ్వు ఇచ్చి \VUgras{శుభ సాయంత్రం} \VUgras{చెప్పేందుకు} వస్తున్న ముద్దుల చిన్నారి
తల్లి. ఆ యువతికి చాలా అందమైన \VUgras{చెవిపోగులు}, హుందా అయిన \VUgras{తలపాగా} ఉన్నాయి.

%%K t04-c1-16-1 p
16. --- తన \VUgras{చొక్కా} \textsuperscript{(36)} వల్లా, ఉబ్బిన \VUgras{చొక్కా చేతుల}
\textsuperscript{(37)} వల్లా అంత విచిత్రంగా కనిపించే ఆ చిన్న చిన్నాన్న కొడుకు చాలా జాలి
కలిగిస్తాడు. అతను \VUgras{అనాథ} \textsuperscript{(35)}. చాలా చిన్నప్పుడే అతను తండ్రినీ తల్లినీ
కోల్పోయాడు. అతని మేనమామే అతని \VUgras{సంరక్షకుడు} \textsuperscript{(38)}, ఆ బిడ్డను పెంచేందుకు
ఆయన పెళ్ళి చేసుకోకుండా ఉండిపోయాడు. ఆయన మంచి మనిషి. ఆయనకు కొంచెం \VUgras{బట్టతల}, చాలా దగ్గరి
చూపు; ఆయన \VUgras{ముక్కు కళ్ళజోడు} వాడతారు.

%%K t04-tit-5 sub
\VUsaut{2.4mm}
\VUpk{10.2pt}{40}{తీపి వంటకాలు.}
\VUsaut{3.0mm}

%%K t04-c1-17-1 p
17. --- విందుకు వచ్చినవాళ్ళ వెనుక, \VUgras{బల్లగుడ్డ} పరిచిన ఒక బల్లపై \VUgras{తీపి వంటకాలు}
\textsuperscript{(39)} సిద్ధం చేసి ఉన్నాయి. ఒక పెద్ద గిన్నెలో ఇదిగో పండ్లు —
\VUgras{పీచుపండ్లు} \textsuperscript{(40)}, \VUgras{నాసపత్తులు} \textsuperscript{(41)},
\VUgras{ఆపిల్స్} \textsuperscript{(42)}, \VUgras{ఆప్రికాట్లు} \textsuperscript{(43)} మరియు
\VUgras{ద్రాక్షలు} \textsuperscript{(44)}. దగ్గరలో \VUgras{అనాసపండ్లు} \textsuperscript{(45)},
అందమైన \VUgras{కేకు} \textsuperscript{(46)}, \VUgras{మద్యం సీసాలు} \textsuperscript{(48)}
మరియు \VUgras{షాంపేన్} \textsuperscript{(49)}, ఇంకా శుభ్రమైన \VUgras{పళ్ళేల}
\textsuperscript{(47)} గుట్టలు.

%%K t04-c1-18-1 p
18. --- \VUgras{మద్యపు గది నౌకరు} \textsuperscript{(50)} పాత ద్రాక్షసారా \VUgras{సీసాను}
\textsuperscript{(52)} \VUgras{కార్క్ తీసే సాధనంతో} \textsuperscript{(51)} తెరుస్తున్నాడు.
\VUgras{కార్క్ మూత} \textsuperscript{(53)} తీయడం చాలా కష్టంగా కనిపిస్తోంది. ఆ నౌకరు చంకలో
\VUgras{సర్వ్ చేసే తువ్వాలు} \textsuperscript{(54)} ఉంది.

%%K t04-c1-19-1 p
19. --- భోజనం తర్వాత పెద్దమనుషులు ధూమపాన గదికి వెళ్తారు, స్త్రీలు నృత్యసభకు సిద్ధమవుతారు.

%%K t04-tit-6 sub
\VUsaut{5.0mm}
\VUcentre{11.4pt}{0}{\textit{రెండవ రంగం.}}
\VUsaut{1.6mm}
\VUpk{11.4pt}{40}{తాతయ్య పుట్టినరోజు.}
\VUsaut{2.6mm}

%%K t04-tit-7 sub
\VUpk{10.2pt}{40}{సందర్శన.}
\VUsaut{3.0mm}

%%K t04-c2-01-1 p
1. --- పది సంవత్సరాలు గడిచాయి, జాన్ తల్లిదండ్రులు వృద్ధులయ్యారు, తమ ముగ్గురు మనవళ్ళంటే —
ఇద్దరు \VUgras{మనవళ్ళు} \textsuperscript{(2)} మరియు ఒక \VUgras{మనవరాలు} \textsuperscript{(3)}
— వాళ్ళకు ఎంతో ప్రేమ. ఈ రోజు \VUgras{తాతయ్య పుట్టినరోజు} కావడంతో పిల్లలు వాళ్ళకు శుభాకాంక్షలు
చెప్పేందుకు వచ్చారు.

%%K t04-c2-02-1 p
2. --- చిన్న పీటరు ఇంకా చాలా చిన్నవాడు, అతనికి మూడేళ్ళే. \VUgras{పిల్లి} \textsuperscript{(1)}
దగ్గరికి రావడం అతను చూస్తున్నాడు. దాని అందమైన పట్టులాంటి \VUgras{బొచ్చునూ} పొడవాటి
\VUgras{తోకనూ} నిమరాలని అతనికి ఉంది; ఆ చక్కని \VUgras{పాదాల} కింద తనను గీరే పదునైన కొంటె
\VUgras{గోళ్ళు} దాగి ఉన్నాయని అతను ఆలోచించడం లేదు.

%%K t04-c2-03-1 p
3. --- అతనికి చేయి ఇచ్చిన అతని అక్క గాబ్రియేలా చాలా గంభీరంగా ఉంది, మార్సెల్లస్ తన పెద్ద
\VUgras{ఓవర్ కోటుతో} \textsuperscript{(4)}, తోలు \VUgras{నడికట్టుతో} \textsuperscript{(5)}
కొంచెం పెద్దవాడిలా కనిపిస్తున్నాడు — తన \VUgras{మేజోళ్ళు} \textsuperscript{(6)} కనిపించేలా
చేసే పొట్టి ప్యాంటు ఉన్నా కూడా.

%%K t04-c2-04-1 p
4. --- వాళ్ళ తండ్రి తన మందపాటి చలికాలపు \VUgras{పెద్ద కోటు} \textsuperscript{(7)}
వేసుకున్నాడు, దానికి \VUgras{బొచ్చు కాలరు} \textsuperscript{(8)} ఉంది; తన \VUgras{జేబులో}
\textsuperscript{(9)} ఒక కర్చీఫ్ ఉంది. అతని భార్యకు \VUgras{ఈకలతో} \textsuperscript{(10)}
అలంకరించిన ఉన్ని టోపీ, తన \VUgras{జాకెట్టు} \textsuperscript{(11)}, తన \VUgras{బొచ్చు కంఠహారం}
\textsuperscript{(12)} మరియు తన \VUgras{ముఫ్} \textsuperscript{(13)} ఉన్నాయి.

%%K t04-c2-05-1 p
5. --- చాలా చలిగా ఉంది, ఎందుకంటే చలికాలం నడుస్తోంది. రోజంతా \VUgras{మంచు}
\textsuperscript{(19)} కురిసింది, ఇళ్ళ కప్పులు పూర్తిగా తెల్లగా ఉన్నాయి. \VUgras{రాత్రితో}
పాటే చలి ఇంకా పదునెక్కుతోంది. ఆకాశంలో \VUgras{చంద్రుడు} \textsuperscript{(17)} మరియు
\VUgras{నక్షత్రాలు} \textsuperscript{(18)} ప్రకాశిస్తూ \VUgras{చీకటిని} చీల్చుకుంటున్నాయి.

%%K t04-tit-8 sub
\VUsaut{4.2mm}
\VUpk{10.2pt}{40}{అనారోగ్యం.}
\VUsaut{3.0mm}

%%K t04-c2-06-1 p
6. --- ఆ పాపం \VUgras{గుడ్డివాడు} \textsuperscript{(20)} చాలా దురదృష్టవంతుడు; అతను
\VUgras{మందుల దుకాణం} \textsuperscript{(16)} ముందున్న చౌక దాటుతున్నాడు, అతన్ని తన
\VUgras{పూడిల్ కుక్క} \textsuperscript{(21)} నడిపిస్తోంది. \VUgras{పనిమనిషి}
\textsuperscript{(14)} యుస్తీనియా వంటగదిలోంచి ఒక \VUgras{గిన్నెడు} \textsuperscript{(15)} చాలా
వేడి \VUgras{కషాయం} తెస్తోంది, దానిలో పంచదార బాగా వేసి ఉంది.

%%K t04-c2-07-1 p
7. --- \VUgras{నాయనమ్మ} \textsuperscript{(23)} బల్లపై తన \VUgras{చెవి కళ్ళజోడు}
\textsuperscript{(26)} వెతకాలనుకుంటోంది, \VUgras{మందుల చీటీ} \textsuperscript{(27)} చదివేందుకు
— ఆ చీటీని \VUgras{వైద్యుడు} \textsuperscript{(25)} ఇప్పుడే రాశాడు; తన \VUgras{కర్రపై}
\textsuperscript{(24)} ఆనుకుని ఆమె తన ఆరామకుర్చీలోంచి లేస్తోంది.

%%K t04-c2-08-1 p
8. --- ఆమె తరచూ \VUgras{చిన్న అనారోగ్యాలతో}, \VUgras{కళ్ళు తిరగడంతో}, \VUgras{జలుబుతో},
\VUgras{పంటినొప్పితో} బాధపడుతుంది; ఆమె కాళ్ళు బలహీనం, ఆమె కళ్ళు ఇక అంత బాగా లేవు. అయినా తనకు
ఎప్పుడూ \VUgras{మందు} \textsuperscript{(28)} రాయాలనుకునే తన ముసలి \VUgras{వైద్యుడిని} ఆమె
ఇష్టంగా ఏడిపిస్తుంది. ఈ రోజు ఆమె చుట్టూ తన యువ కుటుంబం ఉండటంతో ఆమె చాలా సంతోషంగా ఉంది.

%%K t04-c2-09-1 p
9. --- బాగా మూసిన ఆ గదిలో హాయిగా ఉంది. పాలరాతి \VUgras{పొయ్యిలో} \textsuperscript{(29)}
\VUgras{చక్కని మంట} \textsuperscript{(30)} మండుతోంది, ఇప్పుడే వెలిగించిన \VUgras{దీపం}
\textsuperscript{(31)} \VUgras{వెలుగులో} హాయిగా అనిపిస్తోంది.

%%K t04-tit-9 sub
\VUsaut{4.2mm}
\VUpk{10.2pt}{40}{తాతయ్య.}
\VUsaut{3.0mm}

%%K t04-c2-10-1 p
10. --- \VUgras{తాతయ్య} \textsuperscript{(33)} కూడా మరిచిపోయాడు — తాను రోగినని,
\VUgras{కీళ్ళనొప్పులు} తనను \VUgras{ఆరామకుర్చీకి} \textsuperscript{(34)} కట్టిపడేశాయని,
నిన్నటికి నిన్న కూడా తనకు \VUgras{జ్వరమూ సొమ్మసిల్లడమూ} ఉన్నాయని. గత చలికాలంలో తాను
\VUgras{న్యుమోనియాతో} బాధపడ్డాననీ, గరుకైన బాధాకరమైన \VUgras{దగ్గు} తనను చాలా ఇబ్బంది
పెట్టిందనీ అతనికి ఇక గుర్తు లేదు.

%%K t04-c2-10-2 p
అయినా అతను కోలుకోవడం చాలా కష్టమైంది. ఇప్పుడు మాత్రం అతను కాస్త బాగానే ఉన్నాడు. కానీ అతను
\VUgras{దిండుపై} \textsuperscript{(38)} ఆనించిన కాలు కూడా నొప్పి పెడుతోంది — ఆ దిండు ఒక చిన్న
\VUgras{పీటపై} \textsuperscript{(39)} ఉంది.

%%K t04-c2-11-1 p
11. --- తన వెచ్చని \VUgras{ఇంటి అంగీలో} \textsuperscript{(35)} జాగ్రత్తగా చుట్టుకుని
ఉన్నప్పుడు — దానికి ముత్యపు చిప్ప \VUgras{గుండీలు} \textsuperscript{(36)} ఉన్నాయి —, తనకు
వార్తాపత్రిక చదివి వినిపించేందుకు దగ్గర ఆ చిన్న \VUgras{అనాథ బాలిక} \textsuperscript{(40)}
ఉన్నప్పుడు — ఆమె తెల్ల తలపాగా, నీలి \VUgras{గౌను} \textsuperscript{(42)} మరియు
\VUgras{పైవస్త్రం} \textsuperscript{(41)} వేసుకుని ఉంటుంది — అతను ఫిర్యాదు చేయడు.

%%K t04-c2-12-1 p
12. --- ప్రతి సంవత్సరం \VUgras{దినదర్శిని} \textsuperscript{(43)} మళ్ళీ \VUgras{డిసెంబరు}
మొదటి \VUgras{తేదీని} చూపినప్పుడు, అతను తన \VUgras{మనవళ్ళ} కోసం అసహనంగా ఎదురుచూస్తాడు,
ఎందుకంటే ఆ ముఖ్యమైన \VUgras{తేదీని} వాళ్ళు మరిచిపోరని అతనికి తెలుసు.

%%K t04-c2-13-1 p
13. --- చాలా పాతకాలంలో తాను స్వయంగా వెళ్ళి, గోడపై \VUgras{చిత్రపటం} \textsuperscript{(44)}
వేలాడుతున్న ఆ మహిమాన్విత సామ్రాజ్య \VUgras{సేనానికి} శుభాకాంక్షలు చెప్పిన రోజులు అతనికి
భావోద్వేగంతో గుర్తొస్తాయి — ఆ సేనానే తన పిల్లలకు \VUgras{ముత్తాత}.

\VUsaut{6.0mm}
\VUfilet{22mm}
